\NeedsTeXFormat{LaTeX2e}
\documentclass[a4paper,12pt,
headsepline,           % Linie zw. Kopfzeile und Text
oneside,               % einseitig
pointlessnumbers,      % keine Punkte nach den letzten Ziffern in Überschriften
bibtotoc,              % LV im IV
%DIV=15,               % Satzspiegel auf 15er Raster, schmalere Ränder   
BCOR15mm               % Bindekorrektur
%,draft
]{scrbook}
\KOMAoptions{DIV=last} % Neuberechnung Satzspiegel nach Laden von Paket helvet

\pagestyle{headings}
\usepackage{blindtext}

% für Texte in deutscher Sprache
\usepackage[english]{babel}
\usepackage[utf8]{inputenc}
\usepackage[T1]{fontenc}

% Helvetica als Standard-Dokumentschrift
\usepackage[scaled]{helvet}
\renewcommand{\familydefault}{\sfdefault} 

\usepackage{graphicx}

% Literaturverzeichnis mit BibLaTeX
\usepackage[babel,german=quotes]{csquotes}
\usepackage[backend=bibtex8]{biblatex}
\bibliography{bibliography}

% Alternative mit Paket-Option backend=biber und \addbibresource
% \usepackage[backend=biber]{biblatex}
% \addbibresource{bibliography.bib}

% Für Tabellen mit fester Gesamtbreite und variabler Spaltenbreite
\usepackage{tabularx} 

% Besondere Schriftauszeichnungen
\usepackage{url}              % \url{http://...} in Schreibmaschinenschrift
\usepackage{color}            % zum Setzen farbigen Textes

\usepackage{amssymb, amsmath} % Pakete für Mathe-Umgebungen und -Symbole

\usepackage{setspace}         % Paket für div. Abstände, z.B. ZA
%\onehalfspacing              % nur dann, wenn gefordert; ist sehr groß!!
\setlength{\parindent}{0pt}   % kein linker Einzug der ersten Absatzzeile
\setlength{\parskip}{1.4ex plus 0.35ex minus 0.3ex} % Absatzabstand, leicht variabel

% Tiefe, bis zu der Überschriften in das Inhaltsverzeichnis kommen
\setcounter{tocdepth}{3}      % ist Standard

% Beispiele für Quellcode
\usepackage{listings}
\lstset{language=Java,
  showstringspaces=false,
  frame=single,
  numbers=left,
  basicstyle=\ttfamily,
  numberstyle=\tiny}

% hier Namen etc. einsetzen
\newcommand{\fullname}{Christoph Meyer}
\newcommand{\email}{christoph-1.meyer@uni-ulm.de}
\newcommand{\titel}{What people think about Mood Tracker, a systematic review}
\newcommand{\jahr}{2022}
\newcommand{\matnr}{933728}
\newcommand{\gutachterA}{Prof.\,Dr.\,Manfred Reichert}
\newcommand{\betreuer}{Dr.\,Felix Beierle}

% hier die Fakultät auswählen
%\newcommand{\fakultaet}{---  Im Quellcode anpassen nicht vergessen! ---}
\newcommand{\fakultaet}{Ingenieurwissenschaften, Informatik und\\Psychologie}
%\newcommand{\fakultaet}{Mathematik und\\Wirtschafts-\\wissenschaften}
%\newcommand{\fakultaet}{Medizin}
%\newcommand{\fakultaet}{Naturwissenschaften}

% hier das Institut einsetzen
\newcommand{\institut}{Institut für Datenbanken und Informationssysteme (DBIS)}

% Informationen, die LaTeX in die PDF-Datei schreibt
\pdfinfo{
  /Author (\fullname)
  /Title (\titel)
  /Producer     (pdfeTex 3.14159-1.30.6-2.2)
  /Keywords ()
}

\usepackage{hyperref}
\hypersetup{
pdftitle=\titel,
pdfauthor=\fullname,
pdfsubject={Diplomarbeit},
pdfproducer={pdfeTex 3.14159-1.30.6-2.2},
colorlinks=false,
pdfborder=0 0 0	% keine Box um die Links!
}

% Trennungsregeln
\hyphenation{Sil-ben-trenn-ung}

\begin{document}
\frontmatter

% Titelseite
\thispagestyle{empty}
\begin{addmargin*}[4mm]{-10mm}

\hfill
\includegraphics[height=1.8cm]{images/logo_uulm_sw.png}\\[1em]

{\footnotesize
%{\bfseries Universität Ulm} \textbar ~89069 Ulm \textbar ~Germany
\hspace*{115mm}\parbox[t]{35mm}{\bfseries Fakultät für\\
\fakultaet\\
% TODO hier Institut anpassen
\mdseries \institut}\\[2cm]

\parbox{140mm}{\bfseries \LARGE \titel}\\[2.5em]
{\footnotesize Master Thesis Software Engineering}\\[3em]

{\footnotesize \bfseries Vorgelegt von:}\\
{\footnotesize \fullname\\ \email}\\ \matnr\\[2em]
{\footnotesize \bfseries Gutachter:}\\                     
{\footnotesize \gutachterA}\\[2em]
{\footnotesize \bfseries Betreuer:}\\ 
{\footnotesize \betreuer}\\\\
{\footnotesize \jahr}
}
\end{addmargin*}


% Impressum
\clearpage
\thispagestyle{empty}
{ \small
  \flushleft
  Fassung \today \\\vfill
  \copyright~\jahr~\fullname\\[0.5em]
% Wenn Sie Ihre Arbeit unter einer freien Lizenz bereitstellen möchten, können Sie die nächste Zeile in Ihren Code aufnehmen. Bitte beachten Sie, dass Sie hierfür an allen Inhalten, inklusive enthaltener Abbildungen, die notwendigen Rechte benötigen! Beim Veröffentlichungsexemplar Ihrer Dissertation achten Sie bitte darauf, dass der Lizenztext nicht den Angaben in den Metadaten der genutzten Publikationsplattform widerspricht. Nähere Information zu den Creative Commons Lizenzen erhalten Sie hier: https://creativecommons.org/licenses/
%This work is licensed under the Creative Commons Attribution 4.0 International (CC BY 4.0) License. To view a copy of this license, visit \href{https://creativecommons.org/licenses/by/4.0/}{https://creativecommons.org/licenses/by/4.0/} or send a letter to Creative Commons, 543 Howard Street, 5th Floor, San Francisco, California, 94105, USA. \\
  Satz: PDF-\LaTeXe
}

% ab hier Zeilenabstand etwas größer 
\setstretch{1.2}

\tableofcontents

\mainmatter
\chapter{Motivation}
% Was sind moodtracker (defintion, features etc), warum sind diese sinnvoll (Behandlung), Warum
The world health organization (WHO) showed that 20 percent of children and adolescents suffer from mental health conditions. Suicide is the second most death reason among 19-29 years olds\cite{who}.
With shortages in the number of therapists, there is a need for computer-assisted psychotherapy (CAT). 
Self reports and retrospective of habits and feelings are a fundamental concept in improving mental health.
For therapist it's the only chance to get an closer insight in the patients behavior during real live scenarios which helps them to find good strategies and treatments.
The problem with descriptions of events, feelings or behavior that lie in the past is that they usually do not correspond to reality and vary greatly. 
Memories are changed by external circumstances. Thus, negative experiences are perceived more strongly than positive ones\cite{bradburn1987answering}.
Ecological Momentary Assessment (EMA) or diary studies address this issue\cite{shiffman2008ecological}.
In this type of reports, the patient describe their habits on a daily basis, which results in a less biased retrospective because the events are still present.
A additional benefit, Daily journal writing helps patients, to train their own mindfulness and align their own focus on the progress they make. For example Hirage et al. showed that writing a diary can help people after a surgery to set and achieve own goals in their treatment\cite{hiraga2019effects}.
Even for healthy people, journaling can help reduce anxiety and stress and reduce the risk of mental illness.
Designing a text offers a cognitive difficulty, making it more difficult to access and integrate into daily routines.
Mood Tracker apps address this issue by providing a more easy way to track their emotions on their smartphone. 
Online Therapeutic tools like "Moodscope" showed scientific proven improvements of the users mental health\cite{drake2013assessing}. Therefore, mood trackers are serious category in CAT and 14.2 \% of all mobile health applications are mood tracker.
Although the increasing number of mood trackers is a good thing, the quality of these apps varies greatly. 
Scientific reviews of these criticize that many apps are developed without the instruction of psychological professionals and are more in line with the opinions and wishes of users\cite{caldeira_mobile_nodate}\cite{schueller2021understanding}.
But what is the opinion of the users? 
What exactly do users expect from this type of application and how is it implemented? 
While earlier attempt use randomly selected user reviews\cite{caldeira_mobile_nodate}, user interviews\cite{schueller2021understanding} or the mobile Mobile Application Rating Scale (MARS)\cite{myers2020evaluating} we want to use an natural language approach to cover those questions. 
We want to know which issues does user have with those applications and if those issues are in common with earlier researches. Also we want to cover, if academic designs for record mood behavior are in common with user practice. 
Is it more important to track the emotions precisely or is an easy and more accessible representation such as Emojis good enough for user satisfactions?






%Warum braucht es eine neue Studie?
%Main problem with current studies is that they only used randomized samples of the user reviews for predicting their need. Also
% interviews where held which suffer from bad generalizable therefor we are using an NLP based method for analyzing them
% easier to reproduce. 

\chapter{Research Question}
\section{App features}
Those questions are resulting in the feature analysis of the app itself.

\textbf{RQ\_AF01: What kind of representation is used to measure Emotions?}

As Caldeira showed in their research those apps offers different ways to measure emotions\cite{caldeira_mobile_nodate}.
We want to collect the different kinds of representations and look what trend is used to measure them.


\textbf{RQ\_AF02: What data other than emotions can be collected?}

Most of the apps provide additional tracking options, for example sleep conditions.
We want to know what categories are being tracked by the apps beside the emotions.

\textbf{RQ\_AF03: What type of methods are used to motivate users to report their feelings on a daily basis?}

We want to know if there are any patterns to motivate users to use the app on a daily basis.
Are there any reward systems or do they rely more on notifications? 

\textbf{RQ\_AF04: What kind of visualization is used to see review the tracked moods?}
Also the reports are showed differently, you can display your emotions inside a calendar or with a line-chart.
Another option is to see some distribution of your emotions over the last days.
We want to measure, which representations are used within the applications?

\section{User Reviews}
Those questions are resulting on the analyse of user reviews.

\textbf{RQ\_UR01: What are the main reasons in general to like those apps?}

We want to find the main reasons users review for giving those apps a good rating (>4 stars).
Are there any patterns for which does all those apps have in common?

\textbf{RQ\_UR02: What are the main reason in general to dislike those apps?}

What are the reasons in general for dislike those apps? 
Based on review keywords, we want to analyze which topics are the most present in user reviews.

\textbf{RQ\_UR03: How important are different representations of emotions?}

We want to know if other representations than likertscales between good and bad are preferred by users.
While emotions are multidimensional, a limitation to one dimension could be an argument against those applications.
It also categorized the emotions in "good emotions" and "bad emotions" which can be in contrast to the acceptance of negative valences emotions.
In academic research, multidimensional metrics such as the Self-Assessment Manikin\cite{bradley1994measuring} are the preferred choice, but what is the opinion in non-academic self reports?
Based on user reviews, we want to determine how important such metrics are in the functioning of the app.


\textbf{RQ\_UR04: How important are goals for users?}

Caldeira et al. criticized in their review that few apps only track the state and can't be used for setting goals\cite{caldeira_mobile_nodate}.
Those self defined aims can help to work on own patterns and was one of the examine user wishes to improve those apps.
On the one hand, we want to examine whether this trend has changed and on the other see how important this feature actually is for users.  

\textbf{RQ\_UR05: How important is privacy for users of these apps?}

Because of the personal data, privacy concerns was another user issue Caldeira found.
However different studies showed that there is a discrepancy between privacy concerns and the actual behavior\cite{BARTH20171038}.
We want to know if online availability of data is more important than secure the data only on local storage.

\textbf{Has the interest in Mood Tracker increased during the Corona crisis? (optional)}

Social Isolation and loneliness was one of the big challenges during the pandemic\cite{loneliness}.
We want to now if this has an effect in the interest in mood tracker as a type of social support.


\chapter{Related Work}
\section{History of Mood Tracking Apps}
\section{Earlier Reviews of Mood Tracking}
\section{Topic Modeling}
\chapter{Methods}
\section{Data collection}
This study is divided in two separated parts. The first part is a systematic review of the applications. We follow the guidelines of the reporting checklist of the Preferred Reporting Items for Systematic Reviews and Meta-Analysis (PRISMA\cite{moher2015preferred}). 
We define mood trackers as:
"Applications for measuring and reporting mood by themselves on a daily basis".
We exclude applications which are:
\begin{itemize}
    \item mainly focus on journaling and not mood tracking (eg. \textit{Daily Diary: Journal with Lock} \cite{simpleinnovation_2022})
    \item doesn't collect mood behavior (eg. journalistic\cite{engelhardt_2020})
    \item doesn't have a feature for self analyzing mood patterns (eg. \textit{Shmoody: Improve Your Mood} \cite{shmoody})
    \item tracking others behavior such as parenting applications or relationship applications (eg. \textit{behavior tracker} \cite{skycap_2020})
    \item only available as a web app and not included in any applications store (eg. \textit{moodtracker.com}\cite{moodtrackercom})
\end{itemize}
As data source we are using the google play store\cite{google} as well as the apple appstore because the two operation systems covers more than 99\% of the worldwide mobile operation system market share.
Our search queries are: ["mood tracker", "mood journal", "mood ema", "emotion tracker"].
For feature extraction we are using the app descriptions as raw data. We include only applications which are available at the time period of our study (November 2022). 
The User Reviews are crawled with appbot\cite{appbot}.

\section{Analyzing User Reviews}
In research question RQ\_UR1 to RQ\_UR5 we are interested about the context of the different reviews.
There are two types of approaches to get a closer look to that.
An basic approach would be to use n-gram frequency of the topic which are discussed. With the given review stars and additional sentiment analysis we can discover if those topics are rated positiv or negative which is important to know to answer Question RQ\_UR03 - RQ\_UR05. 
Instead of n\_grams we could also use an machine learning approach to learn the topics.
By that we define abstract topics to which the different reviews belong.
While having the prior information on the rating or using semantic segmentation we can check if those reviews are meant to be positive or negativ.



\subsection{Topic Modelling vs. Extractive Topic Tagging vs. Predictive Topic Tagging}
% keyword extraction. Possibilities: 1. Extract with tf-idf -> difficult because tfidf clusters 



\subsection{Tf-IDF and LDA}

\subsection{BertTopic}


\section{Evaluation}

% hier weitere Kapitel einbinden

\appendix
% hier Anhänge einbinden
\chapter{Time Line}



\backmatter


\printbibliography

\clearpage
\thispagestyle{empty}

Name: \fullname \hfill Matrikelnummer: \matnr \vspace{2cm}

\minisec{Erklärung}

Ich erkläre, dass ich die Arbeit selbständig verfasst und keine anderen als die angegebenen Quellen und Hilfsmittel verwendet habe.\vspace{2cm}

Ulm, den \dotfill

\hspace{10cm} {\footnotesize \fullname}
\end{document}
